\documentclass[11pt]{article}

\usepackage[margin=1in]{geometry}
\usepackage{amsmath,amssymb}
\usepackage{booktabs}
\usepackage{graphicx}
\usepackage{natbib}
\usepackage{xcolor}
\usepackage{hyperref}

\hypersetup{
  colorlinks=true,
  linkcolor=blue!55!black,
  citecolor=blue!55!black,
  urlcolor=blue!55!black
}

\graphicspath{
  {../../CoFiTok-internal/artifacts/figures/method_overview_2026-07-08/}
  {../../CoFiTok-internal/artifacts/figures/visual_panels_2026-07-08/}
  {../../CoFiTok-internal/artifacts/figures/summary_2026-07-08/}
  {../../CoFiTok-internal/artifacts/figures/final_visual_claim_2026-07-11/}
}

\newcommand{\method}{CoFiTok}
\newcommand{\epshat}{\hat{\epsilon}}
\newcommand{\xzero}{x_0}
\newcommand{\xt}{x_t}
\newcommand{\zk}{z^{(k)}}
\newcommand{\sk}{S_k}

\title{\method: Coarse-to-Fine Denoising Tokens for Pixel-Space Diffusion}
\author{Anonymous Authors}
\date{}

\begin{document}
\maketitle

\begin{abstract}
Pixel-space diffusion models usually predict a full dense noise field with a
single monolithic network output. We introduce \method, a denoising-token
factorization that represents the dense noise prediction as an ordered sum of
negative-noise components. Each token is expanded by a restricted,
condition-free synthesis operator that only sees the current token, so prefix
sums form valid partial noise predictions and produce controllable partial
denoising results in pixel space. On CIFAR-10, Tiny ImageNet-200, and an
ImageNet-1K 64x64 fallback, a tiny \method{} predictor with restricted synthesis
learns ordered components with clean zero-token diagnostics, strong
random/reverse-order sensitivity, and broader token utilization than
endpoint-only factorized training. At the same 20k-step budget, K8 \method{}
nearly matches endpoint-only final clean MSE on Tiny ImageNet-200 and ImageNet-64 HF
while improving prefix path AUC by roughly 28--35x. Current generated-sample
Frechet and official FID metrics are mixed, so the supported claim is
prefix-controllable denoising at comparable endpoint quality rather than
unconditional generation-quality wins.
\end{abstract}

\section{Introduction}

Diffusion models reverse a noising process by repeatedly predicting a dense
noise field. This dense prediction is effective, but it is also monolithic:
there is no direct way to ask how much of the denoising update has been
explained, whether the update has an interpretable coarse-to-fine order, or how
to stop after a valid partial prediction. Ordered visual tokenizers and
diffusion-autoregressive methods suggest that generation can benefit from
structured intermediate representations. The question in this paper is
narrower: can the dense pixel-space noise prediction itself be factorized into
ordered denoising tokens?

\method{} rewrites the predicted noise field as
\begin{equation}
  \epshat^{(K)} = \sum_{k=1}^{K} S_k\left(z^{(k)}\right),
  \label{eq:full-factorization}
\end{equation}
where each $z^{(k)}$ is a denoising token and $S_k$ expands only
that token into one dense component. Prefixes are meaningful by construction:
\begin{align}
  \epshat^{(1:m)} &= \sum_{k=1}^{m} S_k\left(z^{(k)}\right), \\
  \hat{x}^{(m)}_0 &= \frac{\xt - \sigma_t \epshat^{(1:m)}}{\alpha_t}.
  \label{eq:prefix-denoise}
\end{align}

The central design constraint is that the synthesis operator must not become a
decoder. In the default \method{} model, $S_k$ is token-only, bias-free,
shallow, local, and condition-free. It does not receive the noised image,
timestep, class label, text prompt, previous tokens, skip features, or a learned
constant. This restriction forces the token itself to carry the sample-specific
negative-noise information and gives the zero-token test a concrete meaning:
$S_k(0)=0$ by construction.

The current evidence supports a scoped MVP claim. \method{} learns ordered
denoising components on three public image settings: CIFAR-10, Tiny
ImageNet-200, and an ImageNet-1K 64x64 fallback. Compared with an endpoint-only
epsilon baseline trained for the same 20k steps, the K8 light denoise-path
objective nearly matches final clean MSE on Tiny ImageNet-200 and ImageNet-64
HF, while dramatically improving prefix/path alignment and effective token
usage. Generated-sample distribution smoke metrics do not show a stable
unconditional generation-quality win, which clarifies the present contribution:
\method{} makes dense noise prediction prefix-controllable without relying on a
final VAE-style decoder.

\paragraph{Contributions.}
\begin{enumerate}
  \item We formulate dense pixel-space diffusion noise prediction as an ordered
  sum of denoising-token components.
  \item We introduce a restricted, condition-free synthesis operator whose
  token-only interface and bias-free construction prevent decoder-style
  shortcuts and make $S_k(0)=0$ testable.
  \item We train \method{} with a denoise-path objective that makes token
  prefixes approximate a partial denoising trajectory rather than only the final
  clean image.
  \item We evaluate prefix controllability, component ordering,
  non-degeneration, endpoint quality, and generated-sample smoke metrics across
  CIFAR-10, Tiny ImageNet-200, and ImageNet-64 HF.
\end{enumerate}

\section{Related Work}

\paragraph{Diffusion models and latent diffusion.}
Denoising diffusion probabilistic models learn to reverse a Gaussian noising
process by predicting noise targets \citep{ho2020denoising}. Latent diffusion
models move this process into a learned latent space and decode the final latent
through a VAE-style decoder \citep{rombach2022latentdiffusion}. \method{}
instead stays in the pixel-space denoising algebra and factorizes the dense
noise prediction itself. The \method{} token is not a final image latent; it is a
negative-noise component that directly contributes to the pixel-space
denoising update.

\paragraph{Trajectory ordering and latent-pixel interaction.}
Latent Forcing studies how changing the schedule between latent and pixel
variables changes the information revealed by the diffusion trajectory
\citep{baade2026latentforcing}. \method{} is inspired by this view of trajectory
ordering, but tokenizes the noise prediction inside the reverse trajectory
rather than introducing a separate latent-pixel trajectory.

\paragraph{Ordered visual tokenizers.}
Spectral Image Tokenizer, FlexTok, Matryoshka Representation Learning, M3, and
related nested-token approaches organize image or multimodal representations so
that prefixes or subsets can be useful
\citep{esteves2025spectral,bachmann2025flextok,kusupati2024matryoshka,cai2024matryoshkamultimodal}.
These methods operate on image representations or visual-token budgets.
\method{} operates on dense diffusion noise prediction. Its prefixes are valid
pixel-space denoising updates, not partial decodes from a general image
tokenizer.

\paragraph{Diffusion-autoregressive visual tokens.}
Selftok and D-AR are the closest conceptual neighbors because they connect
diffusion processes and autoregressive tokens \citep{wang2025selftok,gao2025dar}.
\method{} differs in the represented object and in the anti-degeneration
mechanism. \method{} tokens are continuous negative-noise components,
and the synthesis operator is deliberately weak and condition-free. The main
claim is not that diffusion steps are visual tokens, but that dense noise
prediction can be decomposed into ordered components whose prefixes are directly
usable in the diffusion denoising formula.

\section{Method}

\subsection{Forward Noising and Dense Prediction}

For a clean image $x_0$, timestep $t$, and Gaussian noise $\epsilon$, the forward
process is
\begin{equation}
  x_t = \alpha_t x_0 + \sigma_t \epsilon.
  \label{eq:forward}
\end{equation}
A standard epsilon-prediction diffusion model learns
$\epsilon_\theta(x_t,t,c)$. \method{} instead predicts the ordered component sum
in Eq.~\ref{eq:full-factorization}. The current implementation uses
unconditional models for the main MVP evidence. Class labels may exist in
datasets, but the restricted synthesis operator never receives labels, prompts,
timesteps, previous tokens, or image features. Figure~\ref{fig:method} shows the
dataflow and marks the blocked paths into $S_k$ that define the default method
boundary.

\begin{figure}[t]
  \centering
  \includegraphics[width=\linewidth]{method_overview.png}
  \caption{%
  \method{} factorizes dense pixel-space diffusion noise prediction into ordered
  denoising tokens. The predictor may be expressive, but each
  synthesis operator $S_k$ receives only the current token and is restricted to
  shallow, condition-free, bias-free local synthesis. Any prefix of components
  forms a valid partial noise prediction and therefore a valid pixel-space
  denoising result. Dashed blocked paths mark information that must not enter
  $S_k$, including $x_t$, timesteps, labels, prompts, previous tokens, skip
  features, and learned constants.}
  \label{fig:method}
\end{figure}

\subsection{Next Denoising Tokens}

The token predictor $T_k$ may be expressive. It sees the noised image, timestep,
and previous tokens:
\begin{align}
  z^{(1)} &= T_1(x_t,t), \\
  z^{(k)} &= T_k(x_t,z^{(<k)},t), \quad k>1.
\end{align}
In the tiny MVP model, $T_k$ is a small U-Net-like predictor with token heads.
The simultaneous-prediction ablation disables token-to-token feedback so every
token head reads the same hidden state.

\subsection{Restricted Synthesis Operator}

Each token is expanded independently:
\begin{equation}
  \epsilon^{(k)} = S_k\left(z^{(k)}\right).
\end{equation}
The default $S_k$ is a bias-free projection followed by a bias-free local
convolution. It is token-only and condition-free. Because no layer has bias and
there is no learned constant input, zero tokens map to zero components:
\begin{equation}
  S_k(0)=0.
\end{equation}
This restriction separates \method{} from a learned decoder. A deep-$S_k$
ablation is included only to test degeneration risk: it uses biased nonlinear
convolutions while preserving the token-only interface. This ablation can
improve apparent endpoint metrics but fails the zero-token diagnostic, so it is
not the default method.

\subsection{Prefix Denoising}

For any prefix length $m$, \method{} forms the prefix prediction and pixel-space
denoising result in Eq.~\ref{eq:prefix-denoise}. Thus every prefix is a valid
partial noise prediction and every prefix produces a pixel-space denoising
result. This is the primary interface exposed by the method.

\subsection{Denoise-Path Objective}

The current MVP objective is
\begin{equation}
  \mathcal{L}
  =
  \mathcal{L}_{\epsilon}
  + 0.15\,\mathcal{L}_{\mathrm{path\ prefix}}
  + 0.30\,\mathcal{L}_{\mathrm{path\ component}}.
\end{equation}
The denoise path interpolates from the zero-update $x_t$ reconstruction toward
coarse-to-fine spatial targets. The progress schedule uses power 1.5 in the
current best MVP configuration. This objective deliberately avoids forcing early
prefixes to be immediately clean; instead, prefixes are trained to follow a
partial denoising path.

\section{Experiments}

\subsection{Experimental Setup}

\paragraph{Datasets.}
We evaluate on CIFAR-10 as a smoke and diagnostic dataset, Tiny ImageNet-200 as
the main MVP dataset, and an ImageNet-1K 64x64 Hugging Face fallback as an
ImageNet-family validation setting. The preferred TFDS
\texttt{downsampled\_imagenet/64x64} source was attempted but was not available
in the current environment; therefore the fallback is always named explicitly.

\paragraph{Models.}
All current experiments use tiny models. The main method is K8 light
denoise-path \method{} with restricted synthesis. We also evaluate K4 and K16 on
ImageNet-64 HF, epsilon-only endpoint training, channel-mask controls,
simultaneous-predictor ablations, and deep-$S_k$ synthesis ablations.

\paragraph{Metrics.}
We report final clean MSE, prefix denoise-path AUC, effective token count,
zero-token component energy ratio, shuffled-token relative mismatch, PSNR,
LPIPS Alex, low-res Frechet proxy, torchvision Inception Frechet, and
generated-sample smoke-to-50k streamed Frechet. Generated-sample Frechet is
computed up to 8192 generated samples against 8192 real validation images for
the main Tiny/HF rows, up to 50000/50000 for the HF ImageNet-64 fallback, and
with matched 20k multiscale controls on Tiny and HF.
This is still an internal torchvision Inception-Frechet protocol rather than an
external official FID benchmark claim.

\paragraph{Validation.}
The current validation suite passes on the \texttt{pro6000} server. Summary
artifacts contain 111 training reports, 95 order-evaluation reports, 58
fixed-timestep quality reports, 22 sampling reports, and 29 generated-quality
reports. A deterministic full-validation reconstruction and prefix sweep covers
the main Tiny ImageNet-200 10k-validation and ImageNet-64 HF 50k-validation
rows. The stronger multiscale predictor also has matched 20k train-scale rows
and generated-quality controls on Tiny ImageNet-200 and ImageNet-64 HF while
preserving the restricted $S_k$ contract.

\subsection{Endpoint Quality and Prefixes}

At the same 20k training budget, K8 light \method{} nearly matches epsilon-only
final clean MSE while greatly improving path alignment.

\begin{table}[t]
  \centering
  \caption{Matched 20k two-seed comparison (mean $\pm$ sample standard
  deviation over seeds 103 and 139). Lower is better for final clean MSE and
  path AUC; higher is better for effective token count.}
  \label{tab:fair20k}
  \small
  \resizebox{\linewidth}{!}{%
  \begin{tabular}{llrrrr}
    \toprule
    Dataset & Method & Final MSE $\downarrow$ & Path AUC $\downarrow$ & Eff. K $\uparrow$ & Zero ratio $\downarrow$ \\
    \midrule
    \input{../../CoFiTok-internal/artifacts/reports/long_budget_repeat_2026-07-11/long_budget_repeat_rows.tex}\\
    \bottomrule
  \end{tabular}
  }
\end{table}

Across the four paired dataset-seed comparisons, \method{} lowers path AUC in
4/4 cases, with a mean relative reduction of 96.31\%. This structural gain has
a measurable endpoint tradeoff: final MSE increases by 2.89\% on average and
is not lower in any of the four pairs. Figure~\ref{fig:final-visual} shows the
corresponding qualitative prefix grids and synthesis diagnostics.
In the full-validation deterministic sweep, K8 light also improves MSE AUC and
low-res Frechet proxy AUC on both Tiny ImageNet-200 and ImageNet-64 HF, while
endpoint-only factorization keeps the lower endpoint MSE.

\begin{figure}[t]
  \centering
  \begin{minipage}[t]{0.49\linewidth}
    \centering
    \includegraphics[width=\linewidth]{figure2a_tiny_prefix.png}
    \par\smallskip
    {\footnotesize\textbf{(a)} Tiny-IN: endpoint-only (left), \method{} (right).}
  \end{minipage}\hfill
  \begin{minipage}[t]{0.49\linewidth}
    \centering
    \includegraphics[width=\linewidth]{figure2b_imagenet64_prefix.png}
    \par\smallskip
    {\footnotesize\textbf{(b)} ImageNet-64: endpoint-only (left), \method{} (right).}
  \end{minipage}

  \medskip
  \begin{minipage}[t]{0.49\linewidth}
    \centering
    \includegraphics[width=\linewidth]{figure2c_restricted_synthesis.png}
    \par\smallskip
    {\footnotesize\textbf{(c)} Restricted $S_k$: ordered (left), shuffled (right).}
  \end{minipage}\hfill
  \begin{minipage}[t]{0.49\linewidth}
    \centering
    \includegraphics[width=\linewidth]{figure2d_deep_synthesis.png}
    \par\smallskip
    {\footnotesize\textbf{(d)} Deep $S_k$: ordered (left), shuffled (right).}
  \end{minipage}
  \caption{%
  Qualitative prefix and degeneration evidence assembled from four standalone
  subfigure files. In (a--b), each half shows the same examples across the
  prefix path; \method{} produces more organized intermediate denoising than
  endpoint-only factorization. In (c--d), shuffling breaks sample alignment,
  while deep synthesis can carry structure that the restricted, zero-preserving
  operator cannot.}
  \label{fig:final-visual}
\end{figure}

\subsection{Order Sensitivity}

Order ablations keep the same component sum at the endpoint but change the
prefix trajectory. This makes order evaluation a direct test of whether prefixes
are organized rather than an endpoint-quality test.

\begin{table}[t]
  \centering
  \caption{ImageNet-64 HF order sensitivity at seed 139. The endpoint is
  unchanged by order because all components are still summed, but random and
  reverse orders damage prefix path quality, especially for K8/K16.}
  \label{tab:order}
  \begin{tabular}{rlrrr}
    \toprule
    K & Order & Final MSE $\downarrow$ & Path AUC $\downarrow$ & Zero ratio $\downarrow$ \\
    \midrule
    4 & ordered & 0.1029 & 0.0674 & 0.0000 \\
    4 & random & 0.1029 & 0.0727 & 0.0000 \\
    4 & reverse & 0.1029 & 0.7516 & 0.0000 \\
    8 & ordered & 0.1041 & 0.0635 & 0.0000 \\
    8 & random & 0.1041 & 0.2441 & 0.0000 \\
    8 & reverse & 0.1041 & 0.6947 & 0.0000 \\
    16 & ordered & 0.1021 & 0.0680 & 0.0000 \\
    16 & random & 0.1021 & 0.4680 & 0.0000 \\
    16 & reverse & 0.1021 & 0.7020 & 0.0000 \\
    \bottomrule
  \end{tabular}
\end{table}

K8 gives the best current tradeoff: it uses tokens broadly and shows clear
random/reverse degradation. K4 is a strong cost baseline but has less room for a
rich prefix curriculum. K16 does not improve path AUC at 5k steps and likely
requires longer training or stronger shaping before becoming worthwhile. Seed
103 repeats for K4 and K16 preserve the same pattern: reverse order damages
prefix path quality in both cases, while K16 random order is also strongly worse
than ordered accumulation.

\begin{figure}[t]
  \centering
  \includegraphics[width=0.85\linewidth]{imagenet_hf_order_path_auc.png}
  \caption{%
  Component order affects prefix quality even when the final component sum is
  unchanged. Ordered components preserve the denoise path, while random and
  reverse orders damage prefix alignment.}
  \label{fig:order}
\end{figure}

\subsection{Restricted Synthesis and Degeneration}

The restricted synthesis operator has zero-token component energy ratio 0.0 in
the restricted reports. The deep-$S_k$ ablation improves some endpoint/path
numbers but has a nonzero zero-token ratio: final clean MSE 0.0977, path AUC
0.0366, and zero ratio 0.0522. This result supports the design boundary. A
stronger synthesis operator can carry token-independent priors and improve
apparent metrics, but then it no longer satisfies the default \method{}
restriction. A seed-103 repeat shows the same failure mode with zero-token ratio
0.0450. Figure~\ref{fig:final-visual} visualizes restricted/deep and
shuffled-token panels.

\subsection{Component Separation Pressure}

We treat component decorrelation as a pressure-term ablation rather than a
default objective. The headline \method{} rows use the no-decorrelation light
denoise-path objective because it most directly measures ordered prefix
denoising. On Tiny ImageNet-200, a conservative schedule
(\texttt{weight=0.003}, start step 3000, warmup 1500, full weight at step 4500
of 5000) reduces component correlation on two seeds while keeping the
ordered-prefix penalty small.

\begin{table}[t]
  \centering
  \caption{Component-separation pressure ablation on Tiny ImageNet-200. Lower is
  better for ordered path AUC, mean absolute component cosine, and final MSE;
  higher is better for effective token count. Scheduled decor 0.003 consistently
  reduces component cosine while leaving ordered-prefix behavior intact.}
  \label{tab:decor}
  \begin{tabular}{rllrrr}
    \toprule
    Seed & Variant & Ordered AUC $\downarrow$ & Mean abs cosine $\downarrow$ & Final MSE $\downarrow$ & Eff. K $\uparrow$ \\
    \midrule
    139 & baseline & 0.0766 & 0.7223 & 0.1487 & 6.8043 \\
    139 & decor 0.003 & 0.0804 & 0.6890 & 0.1486 & 6.8316 \\
    103 & baseline & 0.0740 & 0.6967 & 0.1503 & 6.8430 \\
    103 & decor 0.003 & 0.0776 & 0.6644 & 0.1501 & 6.8829 \\
    \bottomrule
  \end{tabular}
\end{table}

Across the two seeds, scheduled decor 0.003 changes ordered path AUC by +0.0036,
mean absolute component cosine by -0.0328, final MSE by -0.0002, and effective
token count by +0.0336. Thus mild separation pressure can make components less
correlated without materially breaking ordered prefix behavior. Stronger decor
weights reduce cosine further but impose larger path-AUC penalties, so we keep
them as stress-test ablations rather than the default \method{} objective.

\subsection{Fixed-Timestep and Generated-Sample Quality}

Fixed-timestep 256-image validation shows that epsilon-only remains a strong
endpoint baseline. K8 light is close on PSNR/LPIPS, but it is not uniformly
better on endpoint metrics.

\begin{table}[t]
  \centering
  \caption{Cross-dataset fixed-timestep quality. These are reconstruction probes,
  not sample FID.}
  \label{tab:quality}
  \begin{tabular}{lllrrr}
    \toprule
    Dataset & Method & Steps & Final PSNR $\uparrow$ & Final Inception $\downarrow$ & Final LPIPS $\downarrow$ \\
    \midrule
    CIFAR-10 & epsilon-only & 3k & 14.390 & 334.069 & 0.1639 \\
    CIFAR-10 & K8 light & 3k & 14.170 & 327.889 & 0.1674 \\
    Tiny ImageNet-200 & epsilon-only & 20k & 14.976 & 344.162 & 0.5479 \\
    Tiny ImageNet-200 & K8 light & 20k & 14.865 & 346.083 & 0.5505 \\
    ImageNet-64 HF & epsilon-only & 20k & 15.743 & 330.643 & 0.5204 \\
    ImageNet-64 HF & K8 light & 20k & 15.558 & 360.795 & 0.5150 \\
    \bottomrule
  \end{tabular}
\end{table}

Generated-sample smoke metrics do not establish a stable unconditional
generation-quality win at the current model scale. Earlier 1024/4096 DDIM-20
slices favor epsilon-only on both Tiny ImageNet-200 and ImageNet-64 HF; the
8192/8192 DDIM-50 streamed protocol slightly favors K8 light on Tiny
ImageNet-200 but still favors epsilon-only on ImageNet-64 HF, and the
50000/50000 HF streamed protocol again favors epsilon-only. The matched 20k
multiscale generated-quality controls also favor epsilon-only on both Tiny
ImageNet-200 and ImageNet-64 HF under the internal streamed protocol, while the
external \texttt{pytorch-fid} image-directory protocol is mixed.

\begin{table}[t]
  \centering
  \small
  \caption{Generated-sample smoke-to-50k streamed metrics. Tiny ImageNet-200 is
  shown at 8192/8192 and with a matched multiscale 10000/10000 control; HF
  ImageNet-64 is shown at 8192/8192, 50000/50000, and matched multiscale
  50000/50000. These are internal streamed DDIM Inception-Frechet measurements,
  not external official FID benchmark claims.}
  \label{tab:generated}
  \begin{tabular}{llrrrr}
    \toprule
    Dataset & Method & Generated & Real & Low-res Frechet $\downarrow$ & Inception Frechet $\downarrow$ \\
    \midrule
    Tiny ImageNet-200 & epsilon-only & 8192 & 8192 & 7.5761 & 272.9955 \\
    Tiny ImageNet-200 & K8 light & 8192 & 8192 & 7.3264 & 271.0260 \\
    ImageNet-64 HF & epsilon-only & 8192 & 8192 & 6.5490 & 240.8919 \\
    ImageNet-64 HF & K8 light & 8192 & 8192 & 6.6993 & 269.2785 \\
    ImageNet-64 HF & epsilon-only & 50000 & 50000 & 6.4496 & 238.0329 \\
    ImageNet-64 HF & K8 light & 50000 & 50000 & 6.6096 & 266.8045 \\
    Tiny ImageNet-200 & eps-only ms & 10000 & 10000 & 5.2057 & 136.7132 \\
    Tiny ImageNet-200 & K8 light ms & 10000 & 10000 & 5.7001 & 141.4577 \\
    ImageNet-64 HF & eps-only ms & 50000 & 50000 & 5.5593 & 134.2443 \\
    ImageNet-64 HF & K8 light ms & 50000 & 50000 & 6.2099 & 136.8376 \\
    \bottomrule
  \end{tabular}
\end{table}

\begin{table}[t]
  \centering
  \small
  \caption{Official \texttt{pytorch-fid} image-directory protocol for matched
  20k multiscale checkpoints. Tiny ImageNet-200 is evaluated at 10000/10000
  images and the ImageNet-64 HF fallback at 50000/50000 images. Lower is
  better. These rows are diagnostic evidence, not a stable generation-quality
  superiority claim.}
  \label{tab:official-fid}
  \begin{tabular}{llrr}
    \toprule
    Dataset & Method & Images & Official FID $\downarrow$ \\
    \midrule
    Tiny ImageNet-200 & epsilon-only & 10000 & 148.841 \\
    Tiny ImageNet-200 & K8 light & 10000 & 152.247 \\
    ImageNet-64 HF & epsilon-only & 50000 & 140.721 \\
    ImageNet-64 HF & K8 light & 50000 & 134.082 \\
    \bottomrule
  \end{tabular}
\end{table}

\section{Discussion}

\method{} suggests that dense pixel-space diffusion noise prediction can be
treated as a sequence of ordered denoising-token components. The restricted
synthesis operator is central: it gives the prefix interface meaning and
prevents the method from becoming a generic latent decoder. The denoise-path
objective is also central because early prefixes should represent partial
denoising rather than clean reconstructions.

The strongest empirical result is the combination of near-matched final clean
MSE and large prefix path-AUC gains at 20k steps. The most important negative
result is that generated-sample Frechet smoke metrics do not show a stable
unconditional generation-quality win. Together, these results define the
current scope of the contribution.

\section{Limitations}

The current models are tiny and trained with MVP budgets. Generated-sample
Frechet reaches 50000 generated vs 50000 real images for the HF fallback, and
the external \texttt{pytorch-fid} image-directory protocol has been run for the
matched multiscale 20k checkpoints. These sample-quality results are mixed and
do not show a robust \method{} sample-quality win. Fixed-timestep
Inception/LPIPS metrics use validation slices rather than full-dataset
reconstruction sweeps. Seed repeats cover the main epsilon/K8 comparison but
not every ablation. The ImageNet-family validation uses an HF ImageNet-1K 64x64
fallback because the preferred \texttt{downsampled\_imagenet/64x64} source was
unavailable in the current environment. The venue-neutral draft still needs
author metadata and final checklist decisions before any camera-ready version.

\section{Conclusion}

\method{} factorizes pixel-space diffusion noise prediction into an ordered
  sequence of denoising tokens. Each token is expanded by a restricted,
condition-free synthesis operator into a dense negative-noise component, and
prefix sums yield valid partial denoising results. Across CIFAR-10, Tiny
ImageNet-200, and ImageNet-64 HF, the current MVP evidence shows that \method{}
can learn ordered, non-degenerate prefix-controllable components while
preserving endpoint quality close to an epsilon-only baseline. The next step is
to test whether the same factorization remains useful for larger backbones,
longer training, and official generated-sample FID toolchains.

\bibliographystyle{plainnat}
\bibliography{../references}

\end{document}
